\par\addvspace{12pt}\Needspace{9\baselineskip}

\section{交点与曲线系}

% ITEM: quadratic-example-056
\solutionitem{quadratic-example-056}{例4.56}
（据2024年北京卷第13题）若直线\(y=k(x-3)\)与双曲线\(x^2/4-y^2=1\)只有一个公共点，求\(k\)的所有可能值，并解释它们与渐近线方向的关系。
\solutionanswer
沿直线考察 \(F(x)=x^2/4-k^2(x-3)^2-1\)。在 \(x=0,3\) 处各有一个平方项消失，故 \(F(0)=-9k^2-1<0\)，\(F(3)=5/4>0\)。若 \(F\) 是二次式，它既取正值又取负值，必有两个不同实根。因此只有二次项消失时才可能恰有一个交点：
\[
    \tfrac14-k^2=0,\qquad k=\pm\tfrac12 .
\]
此时 \(F\) 在两处异号，确为非常数的一次式，恰有一根。这两个斜率正是双曲线的渐近线斜率。




% ITEM: quadratic-example-057
\solutionitem{quadratic-example-057}{例4.57}
（2016年江苏卷第21题选修4-4部分）直线\(l\)的参数方程为\(x=1+t/2,y=\sqrt3t/2\)，椭圆\(C\)的参数方程为\(x=\cos\theta,y=2\sin\theta\)。求直线\(l\)截椭圆所得弦\(AB\)的长。
\solutionanswer
参数方向为单位向量。代入椭圆得\(t+7t^2/16=0\)，根为\(0,-16/7\)，所以弦长就是参数差\(16/7\)。


% ITEM: quadratic-example-058
\solutionitem{quadratic-example-058}{例4.58}
求以\(M=(1,1/2)\)为中点、端点在椭圆\(x^2/9+y^2/4=1\)上的弦所在直线及弦长。若中点改为\(N=(3,1)\)，能否仍得到这样的实弦？解释中点条件与交点存在条件各自的作用。
\solutionanswer
设两端点为 \(M\pm(u,v)\)。将它们代入椭圆方程，两式相减得 \(8u+9v=0\)，相加取平均得
\[
 \frac{u^2}{9}+\frac{v^2}{4}=1-\frac19-\frac1{16}=\frac{119}{144}.
\]
第一式说明半弦方向为 \((9,-8)\)，故弦线为 \(8(x-1)+9(y-1/2)=0\)，即 \(16x+18y=25\)。写 \((u,v)=t(9,-8)\)，第二式给出 \(25t^2=119/144\)，所以弦长为 \(2|t|\sqrt{145}=\sqrt{17255}/30\)。若中点换成 \(N=(3,1)\)，相加所得等式的右端为 \(1-1-1/4<0\)。相减仍给出一个方向，相加却说明这个方向上没有实弦。


% ITEM: quadratic-exercise-35-01
\solutionitem{quadratic-exercise-35-01}{习题4.35（1）}
讨论直线\(y=kx+1\)与抛物线\(y^2=4x\)的交点个数，分别处理\(k=0\)与\(k\ne0\)。
\par
\solutionanswer
取 \(y\) 为未知量。由抛物线有 \(x=y^2/4\)，代入直线即得 \(ky^2-4y+4=0\)。\(k=0\) 时只有 \(y=1\)，交点为 \((1/4,1)\)。\(k\ne0\) 时判别式为 \(16(1-k)\)，故 \(k<1\) 时有两个交点，\(k=1\) 时有一个，\(k>1\) 时没有交点。




% ITEM: quadratic-exercise-35-02
\solutionitem{quadratic-exercise-35-02}{习题4.35（2）}
求直线\(x+y=1\)截椭圆\(x^2/9+y^2/4=1\)所得弦的中点与长度，不分别求端点坐标。
\par
\solutionanswer
线上 \(x=1-y\)，故椭圆方程成为 \(4(1-y)^2+9y^2=36\)，配方得 \((13y-4)^2=432\)。两端纵坐标关于 \(4/13\) 对称，纵坐标之差为 \(24\sqrt3/13\)。弦的斜率为 \(-1\)，故中点为 \((9/13,4/13)\)，弦长为 \(24\sqrt6/13\)。

也可令 \(u=x/3,\ v=y/2\) 化椭圆为单位圆。圆心在弦 \(3u+2v=1\) 上的垂足为 \((3,2)/13\)，还原就是上述中点。圆中半弦长为 \(\sqrt{1-1/13}\)；其单位方向 \((2,-3)/\sqrt{13}\) 还原为 \((6,-6)/\sqrt{13}\)，故原弦长为 \(2\sqrt{1-1/13}\cdot6\sqrt{2/13}=24\sqrt6/13\)。




% ITEM: quadratic-exercise-35-03
\solutionitem{quadratic-exercise-35-03}{习题4.35（3）}
过抛物线\(y^2=4x\)的焦点作弦\(AB\)，其中点为\(M\)。求\(M\)的轨迹及实际范围，说明边界点对应哪条弦。
\par
\solutionanswer
水平线不成焦弦，其余过焦点的直线都可写成 \(x=1+my\)，包括竖直线。与抛物线联立，只对纵坐标配一次方：
\[
    y^2=4(1+my)\quad\Longleftrightarrow\quad (y-2m)^2=4(1+m^2).
\]
右端恒正，所以每个 \(m\) 都给出一条实弦，两端的纵坐标关于 \(2m\) 对称。中点就在弦线上，故 \(y_M=2m,\ x_M=1+my_M=1+y_M^2/2\)。当 \(m\) 遍历实数时 \(y_M\) 也遍历实数，轨迹是整条抛物线；最小横坐标 \(1\) 对应 \(m=0\)，即竖直焦弦 \(x=1\)。


% ITEM: quadratic-exercise-35-04
\solutionitem{quadratic-exercise-35-04}{习题4.35（4）}
求椭圆\(x^2/9+y^2/4=1\)中以\((1,0)\)、\((0,1)\)为中点的弦。若中点为原点，方向是否唯一？
\par
\solutionanswer
令 \(u=x/3,\ v=y/2\) 化椭圆为单位圆。在圆内给定非圆心的弦中点，弦只能垂直于过它的半径。两个给定中点分别落在水平、竖直半径上，还原后所求弦因此分别为 \(x=1\) 和 \(y=1\)。代入原椭圆，前一弦的端点为 \((1,\pm4\sqrt2/3)\)，长 \(8\sqrt2/3\)；后一弦的端点为 \((\pm3\sqrt3/2,1)\)，长 \(3\sqrt3\)。若中点为原点，则过中心的每条直径都合乎要求。




% ITEM: quadratic-exercise-35-05
\solutionitem{quadratic-exercise-35-05}{习题4.35（5）}
一条直线依次与两个同心位似椭圆\(x^2/a^2+y^2/b^2=1\)、\(x^2/a^2+y^2/b^2=4\)相交。证明两条弦具有相同中点，从而直线上位于两椭圆之间的两段长度相等。
\par
\solutionanswer
令 \(u=x/a,\ v=y/b\)，两个椭圆便同时化为同心圆。从圆心向割线作垂线，垂足同时平分两条弦。仿射变换保持中点，还原后两弦仍有同一个中点，故外弦超出内弦的两段等长。

用本节的参数法看，沿割线设 \(X=P+td\)，两椭圆给出的二次方程只有常数项不同，所以两根的平均值相同。这就是同一个中点。




% ITEM: quadratic-exercise-35-06
\solutionitem{quadratic-exercise-35-06}{习题4.35（6）}
设\(d\)是单位向量，将它改为\(2d\)后，式\textup{(4.120)}中的三个系数、根及弦长公式分别如何变化？
\par

\par\smallskip\noindent 正文所引关系为：
\begin{equation*}
\begin{aligned}
    \alpha&=Au^2+Buv+Cv^2,\\
    \beta&=(2Ax_0+By_0+D)u+(Bx_0+2Cy_0+E)v,\\
    \gamma&=\Phi(x_0,y_0).
    \end{aligned}
\end{equation*}
\solutionanswer
同一点原来写成 \(P+td\)，改用 \(2d\) 后就写成 \(P+s(2d)\)，其中 \(t=2s\)。因此只须在原式中换元：
\[
    \alpha t^2+\beta t+\gamma=4\alpha s^2+2\beta s+\gamma .
\]
三个系数变为 \(4\alpha,2\beta,\gamma\)，每个根减半。弦长公式 \(|s_1-s_2|\,|2d|\) 中，参数差减半、方向向量的长度加倍，弦长不变。




% ITEM: quadratic-exercise-35-07
\solutionitem{quadratic-exercise-35-07}{习题4.35（7）}
直线\(x=1+2t,y=t\)截椭圆\(x^2/9+y^2/4=1\)于\(A,B\)。不分别求两个交点，求弦长、中点及\(|PA|\,|PB|\)，其中\(P=(1,0)\)。用有向长度参数核对。
\par
\solutionanswer
弦长用两根之差，中点用两根之和，距离积用两根之积。令 \(X=P+t(2,1)\)，代入椭圆得 \(25t^2+16t-32=0\)。记两根之和为 \(s=-16/25\)，积为 \(p=-32/25\)，则
\[
 M=P+\tfrac{s}{2}(2,1)=(9/25,-8/25),\qquad
 |AB|=\sqrt5\,\sqrt{s^2-4p}=\frac{24\sqrt{30}}{25}.
\]
沿弦每增加一个单位参数，实际前进 \(\sqrt5\)，故 \(|PA|\,|PB|=5|p|=32/5\)。若直接用有向长度 \(l=\sqrt5\,t\)，方程变为 \(5l^2+(16/\sqrt5)l-32=0\)，两根之积为 \(-32/5\)，根差的绝对值也正是上述弦长。


% ITEM: quadratic-exercise-35-08
\solutionitem{quadratic-exercise-35-08}{习题4.35（8）}
（据2015年广东卷理科第20题）过原点的动直线交圆 \(x^2+y^2-6x+5=0\) 于不同两点。求圆心及弦中点的轨迹 \(C\)；若直线 \(y=k(x-4)\) 与 \(C\) 恰有一个交点，求 \(k\) 的范围。
\par
\solutionanswer
圆心 \(K=(3,0)\)，半径为 \(2\)。弦中点 \(M=(x,y)\) 是 \(K\) 在割线上的垂足，故 \(x(x-3)+y^2=0\)，而不同端点要求 \(KM^2=9-3x<4\)。反过来，满足这两个条件的点都是相应弦的中点。轨迹是圆 \(x^2+y^2=3x\) 上 \(x>5/3\) 的弧，含 \((3,0)\)，不含两个端点 \((5/3,\pm2\sqrt5/3)\)。

从 \(T=(4,0)\) 看这段弧。过任一端点的直线斜率绝对值为 \(k_0=2\sqrt5/7\)；从 \(T\) 向轨迹所在的整圆所作切线的斜率绝对值为 \(k_1=3/4\)（圆心到直线的距离为 \(3/2\)），切点横坐标为 \(12/5>5/3\)。当 \(|k|\) 从零增大时，最初只有靠近 \(T\) 的交点在弧上；越过 \(k_0\) 后远交点也进入，到 \(k_1\) 时两点合一。端点本身不取，故恰有一个交点的条件是
\[
    |k|\le 2\sqrt5/7\quad\text{或}\quad |k|=3/4 .
\]


% ITEM: quadratic-exercise-35-09
\solutionitem{quadratic-exercise-35-09}{习题4.35（9）}
（据2024年北京卷第19题）椭圆 \(x^2/a^2+y^2/b^2=1\)（\(a>b>0\)）的两焦点及两短轴端点组成边长为 \(2\) 的正方形。求椭圆方程及离心率。过 \((0,t)\)（\(t>\sqrt2\)）的非竖直割线交椭圆于 \(A,B\)，直线 \(AC\)（\(C=(0,1)\)）再交椭圆于 \(D\)。若 \(BD\) 水平，求 \(t\)。
\par
\solutionanswer
正方形的两条对角线给出 \(b=c=\sqrt2\)，故 \(a=2,\ e=1/\sqrt2\)，椭圆为 \(x^2+2y^2=4\)。

只把横坐标缩为原来的 \(1/\sqrt2\)，椭圆就变成半径 \(b=\sqrt2\) 的圆，纵轴上的 \(C=(0,1)\)、\(T=(0,t)\) 都不动。仍用原字母记各点的像。因 \(BD\) 水平，上、下轴端点 \(U=(0,b),V=(0,-b)\) 分别平分 \(BD\) 的两段圆弧。等弧所对的圆周角相等，故 \(AU,AV\) 是直线 \(AB,AD\) 的两条角平分线。在三角形 \(ACT\) 中，\(U\) 在线段 \(CT\) 内，\(V\) 在其延长线上，由内、外角平分线定理，
\[
    \frac{b-1}{t-b}=\frac{AC}{AT}=\frac{b+1}{t+b}.
\]
解得 \(t=b^2=2\)。

也可留在原坐标中。\(BD\) 水平，故可设 \(B=(w,z),D=(-w,z),A=(u,v)\)。两直线 \(AB,AD\) 在 \(y\) 轴上的截距分别为 \(t,1\)。把两个截距相乘，含 \(uw\) 的交叉项便消去了：
\[
 t\cdot1=\frac{uz-wv}{u-w}\frac{uz+wv}{u+w}
        =\frac{u^2z^2-w^2v^2}{u^2-w^2}=2,
\]
末步只用了 \(u^2=4-2v^2,\ w^2=4-2z^2\)。例如取 \(A=(2,0),D=(-2/3,4/3),B=(2/3,4/3)\)，即得到这样的割线。


% ITEM: quadratic-example-059
\solutionitem{quadratic-example-059}{例4.59}
椭圆\(x^2/4+y^2=1\)与直线\(x+2y=1\)交于\(A,B\)，\(P=(0,1)\)。将原点移到\(P\)，求\(PA,PB\)的斜率之和与积，再用直接代入核对。
\solutionanswer
取\(u=x,v=y-1\)，直线条件为\(u+2v=-1\)，椭圆齐次化后为\(u^2-8uv-12v^2=0\)。两方向斜率和为\(-2/3\)，积为\(-1/12\)。直接代入\(y=(1-x)/2\)得\(2x^2-2x-3=0\)，再对\((y-1)/x\)用根与系数关系，可核对同一结果。


% ITEM: quadratic-example-060
\solutionitem{quadratic-example-060}{例4.60}
设\(P=(0,1)\)，\(A,B\)为椭圆\(x^2/4+y^2=1\)上异于\(P\)的两个不同点，两连线的斜率满足\(k_{PA}k_{PB}=-1/2\)。证明弦\(AB\)经过定点\((0,-1/3)\)。
\solutionanswer
取 \(u=x,\ v=y-1\)，以 \(P\) 为原点。弦不过 \(P\)，设为 \(mu+nv=1\)。在两端点，椭圆方程 \(u^2/4+v^2+2v=0\) 的一次项可写成 \(2v(mu+nv)\)。除以 \(u^2\)，两条连线的斜率 \(r=v/u\) 就是
\[
 (1+2n)r^2+2mr+1/4=0
\]
的两根。根积为 \(-1/2\)，故 \(n=-3/4\)。在弦方程中令 \(u=0\)，无论 \(m\) 为何都有 \(v=-4/3\)，即所有弦都过 \((0,-1/3)\)。

若要先猜这个点，可取关于 \(y\) 轴对称的两个端点：斜率为 \(\pm1/\sqrt2\)，将 \(x^2=2(y-1)^2\) 代入椭圆，除去 \(P\) 就得到水平弦 \(y=-1/3\)。


% ITEM: quadratic-exercise-36-01
\solutionitem{quadratic-exercise-36-01}{习题4.36（1）}
从\(P=(0,1)\)观察椭圆\(x^2/4+y^2=1\)与直线\(y=kx+m\)的两个交点。用移轴齐次化求两方向的斜率和与积，明确所用表达的范围。
\par
\solutionanswer
移轴\(u=x,v=y-1\)，椭圆为\(u^2/4+v^2+2v=0\)，直线为\(ku-v=1-m\)。齐次化得\((1-m)u^2/4+2kuv-(1+m)v^2=0\)。所以两方向斜率和为\(2k/(1+m)\)，积为\((m-1)/(4(m+1))\)。这些有限斜率表达要求\(m\ne\pm1\)，并且原直线有两个不同实交点；\(m=1\)时一个交点就是观察点，\(m=-1\)时一方向竖直，须用方向方程处理。


% ITEM: quadratic-exercise-36-02
\solutionitem{quadratic-exercise-36-02}{习题4.36（2）}
椭圆\(x^2/9+y^2/4=1\)上两点\(A,B\)均异于\(P=(3,0)\)。若\(k_{PA}+k_{PB}=s\)为定值，判断弦\(AB\)的直线系特征，是平行还是过定点，并求相应方向或定点。
\par
\solutionanswer
将原点移到\(P\)，令\(u=x-3,\ v=y\)，椭圆方程变为\(4u^2+9v^2+24u=0\)。弦线\(AB\)不过\(P\)，故可设为\(mu+nv=1\)。在\(A,B\)处有\(24u=24u(mu+nv)\)，代入即得\(PA,PB\)的方向方程
\[
 (4+24m)u^2+24nuv+9v^2=0.
\]
令\(r=v/u\)，则两方向的斜率是\(9r^2+24nr+4+24m=0\)的两根。由根与系数的关系，\(s=-8n/3\)，所以弦线方程化为\(mu-(3s/8)v=1\)，只有\(m\)随弦变化。

若\(s=0\)，弦线为\(mu=1\)，均平行于\(y\)轴。若\(s\ne0\)，令\(u=0\)即得\(v=-8/(3s)\)，故所有弦线都经过原坐标中的定点\((3,-8/(3s))\)。


% ITEM: quadratic-exercise-36-03
\solutionitem{quadratic-exercise-36-03}{习题4.36（3）}
抛物线\(y^2=4x\)与直线\(3x-2y=4\)交于\(A,B\)。先用齐次化求\(OA,OB\)的方向，再求两交点坐标。
\par
\solutionanswer
在直线上用\(4x=x(3x-2y)\)，得到\(3x^2-2xy-y^2=(3x+y)(x-y)=0\)。两方向为\(y=-3x,y=x\)，分别联立原直线得\((4/9,-4/3)\)、\((4,4)\)。


% ITEM: quadratic-exercise-36-04
\solutionitem{quadratic-exercise-36-04}{习题4.36（4）}
圆\(x^2+y^2=4\)与直线\(x+y=1\)交于\(A,B\)，\(P=(1,1)\)。以\(P\)为新原点作齐次化，求\(PA,PB\)的方向方程、斜率和与斜率积。
\par
\solutionanswer
令\(u=x-1,v=y-1\)，则\(u+v=-1\)，圆为\(u^2+v^2+2(u+v)-2=0\)。利用直线条件齐次化得\(3u^2+8uv+3v^2=0\)，故斜率和为\(-8/3\)，积为\(1\)。


% ITEM: quadratic-exercise-36-05
\solutionitem{quadratic-exercise-36-05}{习题4.36（5）}
椭圆\(x^2/a^2+y^2/b^2=1\)（\(a>b>0\)）内，过非中心点\(S=(s,0)\)作两条垂直弦，中点为\(M,N\)。求\(MN\)的定点，再取\(S\)为焦点。提示：先求中点轨迹。
\par
\solutionanswer
以 \(S\) 为原点，令 \(u=x-s,\ v=y\)。若 \(M=(u,v)\) 是过 \(S\) 的一条弦的中点，则半弦方向为 \((u,v)\)（\(M=S\) 的轴向情形最后处理）。由中点弦条件，
\[
    (u+s)u/a^2+v^2/b^2=0 .
\]
反过来，此式上的非零点确定一条过内点 \(S\) 的弦，其中点正是该点。\(N\) 也在这条中点椭圆上。

设 \(MN:\ hu+jv=1\)。在上式的一次项中用 \(su(hu+jv)\) 代替 \(su\)，所得方向方程为
\[
 (1+sh)u^2/a^2+sjuv/a^2+v^2/b^2=0 .
\]
两条弦垂直，故斜率积 \(b^2(1+sh)/a^2=-1\)。于是 \(h=-(a^2+b^2)/(sb^2)\)，而 \(j\) 不必求。令 \(v=0\)，还原得定点 \((a^2s/(a^2+b^2),0)\)。两弦沿坐标轴时，中点是 \(O,S\)，其连线也过此点。取右焦点只须令 \(s=c\)。


% ITEM: quadratic-exercise-36-06
\solutionitem{quadratic-exercise-36-06}{习题4.36（6）}
（据2016年全国II卷文科第21题）\(A\) 是椭圆 \(x^2/4+y^2/3=1\) 的左顶点，过 \(A\) 的两条垂直弦为 \(AM,AN\)，\(AM\) 的斜率为 \(k>0\)。若 \(|AM|=|AN|\)，求 \(\triangle AMN\) 的面积；若 \(2|AM|=|AN|\)，证明 \(\sqrt3<k<2\)。
\par
\solutionanswer
以 \(A=(-2,0)\) 为原点，取 \(u=x+2,\ v=y\)。椭圆化为 \(u^2/4+v^2/3=u\)。弦 \(v=ku\) 的一个端点已是原点，约去 \(u\) 即得另一端的横坐标 \(12/(3+4k^2)\)。所以弦长 \(L(k)=12\sqrt{1+k^2}/(3+4k^2)\)。另一条弦的斜率是 \(-1/k\)，于是
\[
 \frac{|AN|}{|AM|}=k\,\frac{4k^2+3}{3k^2+4}=:f(k).
\]
分式 \((4z+3)/(3z+4)=4/3-7/[3(3z+4)]\) 对 \(z>0\) 递增，故 \(f\) 在正半轴严格递增。\(f(1)=1\)，等长时 \(k=1\)，面积 \(L(1)^2/2=144/49\)。又 \(f(\sqrt3)=15\sqrt3/13<2\)（\(675<676\)），\(f(2)=19/8>2\)，所以 \(f(k)=2\) 时 \(\sqrt3<k<2\)。


% ITEM: quadratic-exercise-36-07
\solutionitem{quadratic-exercise-36-07}{习题4.36（7）}
（据2022年浙江卷第21题）\(A,B\) 是椭圆 \(x^2/12+y^2=1\) 上异于 \(P=(0,1)\) 的两点，\(Q=(0,1/2)\) 在线段 \(AB\) 上。\(PA,PB\) 分别交 \(y=-x/2+3\) 于 \(C,D\)。求 \(P\) 到椭圆上点的最大距离及 \(|CD|\) 的最小值。
\par
\solutionanswer
椭圆上一点到 \(P\) 的距离平方为 \(12(1-y^2)+(y-1)^2=144/11-11(y+1/11)^2\)，故最大距离为 \(12/\sqrt{11}\)。

把原点移到 \(P\)，取 \(u=x,\ v=y-1\)。弦过 \(Q\)，可设 \(hu-2v=1\)。在椭圆方程 \(u^2/12+v^2+2v=0\) 中齐次化，得 \(u^2/12+2huv-3v^2=0\)。因此 \(PA,PB\) 的斜率 \(r_1,r_2\) 之积恒为 \(-1/36\)，只需保留其和 \(s=2h/3\)。直线 \(v=r_i u\) 与 \(v=2-u/2\) 交于横坐标 \(2/(r_i+1/2)\)，故
\[
 |CD|=\sqrt5\,\frac{\sqrt{s^2+1/9}}{|s/2+2/9|}
 \ \ge\ \frac{6\sqrt5}{5}.
\]
末步对 \((s,1/3)\) 与 \((1/2,2/3)\) 用柯西不等式。取 \(s=1/4\)，即 \(h=3/8\)，两向量同向，弦过内点 \(Q\) 且两交点 \(C,D\) 均存在，故下界可达。


% ITEM: quadratic-example-061
\solutionitem{quadratic-example-061}{例4.61}
过\(P=(2,0)\)的两条直线分别交双曲线\(x^2-y^2=1\)于\(A,B\)和\(C,D\)，方向分别为\((1,0)\)、\((2,1)\)。求两组有向长度乘积之比，并说明为何要先将方向单位化。再讨论过\(P=(1,0)\)的直线截抛物线\(y^2=4x\)时，长度乘积如何随方向变化。
\solutionanswer
双曲线中\(\Phi(P)=3\)，两个单位方向的二次型值为\(1,3/5\)，故有向长度乘积分别为\(3,5\)，比为\(3/5\)。若直接用\((2,1)\)的参数，必须再乘模平方\(5\)，不能把参数乘积当长度乘积。对抛物线，沿过焦点的单位方向取点\((1+s\cos\theta,s\sin\theta)\)。代入得\(s^2\sin^2\theta-4s\cos\theta-4=0\)。当\(\sin\theta\ne0\)时，两根异号，有向长度乘积为\(-4/\sin^2\theta\)，无向乘积为\(4/\sin^2\theta\)。水平方向只有顶点这一个有限交点，不构成上述两段长度的乘积。


% ITEM: quadratic-exercise-37-01
\solutionitem{quadratic-exercise-37-01}{习题4.37（1）}
过\(P=(1,0)\)的两条不同直线，分别与抛物线\(y^2=4x\)交于\(A,B\)和\(C,D\)。若两组无向长度乘积相等，证明两条直线关于\(x\)轴对称。
\par
\solutionanswer
焦点\((1,0)\)处弦的距离乘积为\(4/\sin^2\theta\)，其中\(\theta\)为方向角，水平线不是有两个有限端点的弦。乘积相等迫使两方向的\(\sin^2\theta\)相同；两直线不同，只能关于\(x\)轴对称。


% ITEM: quadratic-exercise-37-02
\solutionitem{quadratic-exercise-37-02}{习题4.37（2）}
过\(P=(0,0)\)的直线与双曲线\(x^2/4-y^2=1\)相交。求可作割线的方向范围，并求\(|PA|\,|PB|\)的取值范围。
\par
\solutionanswer
斜率\(k\)须满足\(1/4-k^2>0\)，即\(|k|<1/2\)。单位方向下距离乘积为\((1+k^2)/(1/4-k^2)\)，范围为\([4,+\infty)\)。下端在水平直径取到，接近渐近方向时无上界。


% ITEM: quadratic-exercise-37-03
\solutionitem{quadratic-exercise-37-03}{习题4.37（3）}
对椭圆\(x^2/9+y^2/4=1\)，过内点\(P\)作两条互相垂直的弦\(AB,CD\)。证明\((|PA|\,|PB|)^{-1}+(|PC|\,|PD|)^{-1}\)与方向无关，写出它与\(P\)的关系。
\par
\solutionanswer
单位垂直方向\(d,e\)满足\(q(d)+q(e)=1/9+1/4=13/36\)。分别用方向幂相除，倒数和为\(13/[36(1-q(P))]\)，只依赖内点\(P\)。


% ITEM: quadratic-example-062
\solutionitem{quadratic-example-062}{例4.62}
椭圆\(E:x^2/9+y^2/4=1\)与直线\(l:16x+18y=25\)交于\(A,B\)，\(l\)与\(x=1\)交于\(M\)。过\(M\)作两条不同的弦\(CD,EF\)，直线\(CE,DF\)分别与\(l\)交于\(U,V\)。证明\(x_U+x_V=2\)、\(y_U+y_V=1\)。
\solutionanswer
由 \(x=1\) 得 \(M=(1,1/2)\)。弦线 \(16x+18y=25\) 正是例4.58求出的以 \(M\) 为中点的弦线。由一般蝴蝶定理，\(M\) 也平分 \(UV\)，故 \(U+V=2M=(2,1)\)。

也可用下一小节的曲线束验证。沿 \(l\) 以 \(M\) 为参数原点：椭圆方程 \(F=0\) 的两根互为相反数，故无一次项；弦线 \(CD,EF\) 的方程之积 \(L\) 则是参数平方的倍数。交叉连线 \(CE,DF\) 组成的退化成员 \(\lambda F+\mu L=0\) 沿 \(l\) 也无一次项，故交点 \(U,V\) 的参数互为相反数。


% ITEM: quadratic-exercise-38-01
\solutionitem{quadratic-exercise-38-01}{习题4.38（1）}
抛物线\(y^2=4x\)上取\(A=(1,2)\)、\(B=(1,-2)\)，中点为\(M=(1,0)\)。过\(M\)任取两条不同的非水平弦\(CD,EF\)，直线\(CE,DF\)与\(x=1\)分别交于\(U,V\)。证明\(y_U+y_V=0\)。若改以\(A=(1,2),B=(4,4)\)为起始弦，结论应怎样表述？
\par
\solutionanswer
由一般蝴蝶定理，\(M\)也是\(UV\)的中点，所以\(y_U+y_V=0\)。换起始弦后\(M=(5/2,3)\)，同理\(U+V=2M=(5,6)\)，即两坐标和分别为\(5,6\)。各结论都针对所述交点存在的配置。


% ITEM: quadratic-exercise-38-02
\solutionitem{quadratic-exercise-38-02}{习题4.38（2）}
对双曲线\(x^2/4-y^2=1\)，取竖直弦\(x=3\)及其中点\(M\)。写出适用蝴蝶定理的配置，说明为什么不应将渐近方向当作其中的弦方向。
\par
\solutionanswer
竖直弦端点为\((3,\pm\sqrt5/2)\)，中点\(M=(3,0)\)。过\(M\)选两条各有两个不同有限交点的弦\(CD,EF\)，使\(CE,DF\)与\(x=3\)相交，则两交点之和为\((6,0)\)。渐近方向使限制式降为一次，不能提供两个有限端点，故不能当作这类弦方向。


% ITEM: quadratic-exercise-38-03
\solutionitem{quadratic-exercise-38-03}{习题4.38（3）}
在证明中交换\(E,F\)，证明\(CF,DE\)与\(AB\)的两交点也被\(M\)平分。
\par
\solutionanswer
交换\(E,F\)不改变通过\(M\)的弦\(EF\)及其交点配对条件。对同一定理应用新的标记，就得到\(CF,DE\)与\(AB\)的两交点之和仍为\(2M\)。


% ITEM: quadratic-exercise-38-04
\solutionitem{quadratic-exercise-38-04}{习题4.38（4）}
在蝴蝶定理的证明中，若把\(AB\)所在直线改为任意过\(M\)的直线，方程中通常会重新出现\(y\)的一次项。解释中点条件如何使这一项消失，并将它与4.3的直径方程联系起来。
\par
\solutionanswer
取\(M\)为原点、起始弦为一坐标方向，限制式的两根关于零对称，根和为零，恰使该方向的一次系数消失。改取任意方向时通常不再有此性质。将\(M+td\)代入原二次式并令\(t\)的一次系数为零，就是4.3中与方向\(d\)相应的直径条件。


% ITEM: quadratic-exercise-38-05
\solutionitem{quadratic-exercise-38-05}{习题4.38（5）}
曲线\((y-x)^2=4(x+y)\)上一条弦的中点为\(M=(1,1)\)。求这条弦所在直线，再写出相应蝴蝶配置中\(U,V\)的坐标和。
\par
\solutionanswer
在\(M\pm t(u,v)\)处代入并相减，得到\(u+v=0\)，故弦线\(x+y=2\)。两端点为\((1\pm\sqrt2,1\mp\sqrt2)\)，确为实弦。一般蝴蝶定理给出\(U+V=2M=(2,2)\)。


% ITEM: quadratic-example-063
\solutionitem{quadratic-example-063}{例4.63}
求\(x^2+y^2=5\)与\(x^2+4y^2=8\)的全部交点。能否不求四次方程，而从两个方程的线性组合先找出包含这些点的直线？
\solutionanswer
相减得\(3(y^2-1)=0\)，先得到两条水平线\(y=\pm1\)，再代入圆得\(x=\pm2\)。全部交点为\((\pm2,\pm1)\)。


% ITEM: quadratic-example-064
\solutionitem{quadratic-example-064}{例4.64}
对\(F=x^2+y^2-5\)、\(G=x^2+4y^2-8\)，求曲线束中经过原点的成员，并判断其类型。这个成员能否再与\(F=0\)联立，恢复全部基点？为什么？
\solutionanswer
过原点要求\(-5\lambda-8\mu=0\)，可取成员\(8F-5G=3(x^2-4y^2)=0\)，是两条相交直线。与\(F=0\)联立后必有\(G=0\)，因为其系数\(-5\ne0\)，所以可恢复全部基点。


% ITEM: quadratic-example-065
\solutionitem{quadratic-example-065}{例4.65}
利用曲线束求\(y=x^2\)与\(x=y^2\)的全部实交点，再求束中过\((2,0)\)的成员，并判断类型。
\solutionanswer
两原方程相减给出\((y-x)(1+x+y)=0\)，联立得实交点\((0,0),(1,1)\)，另一直线\(x+y=-1\)不产生实交点。过\((2,0)\)的成员可取\((y-x^2)+2(x-y^2)=0\)，配方为\((x-1)^2+2(y-1/4)^2=9/8\)，是椭圆。


% ITEM: quadratic-example-066
\solutionitem{quadratic-example-066}{例4.66}
过\(P=(1,1)\)作\(y-1=\pm(x-1)\)，与抛物线\(y^2=2x\)交于四点。不求四点坐标，求经过它们的圆。再将两条直线改为\(y-1=\pm k(x-1)\)，求四点互不相同的\(k\)的范围及对应圆系。
\solutionanswer
将两线积记为\(L=(y-1)^2-k^2(x-1)^2\)，抛物线式记为\(F=y^2-2x\)。比较平方项，取\((1+k^2)F-L=0\)，得到圆系
\[
k^2(x^2+y^2)-(2+4k^2)x+2y+k^2-1=0.
\]
\(k=1\)时为\(x^2+y^2-6x+2y=0\)。当 \(k\ne0\) 时，沿任一给定直线，\(F=y^2-2x\) 在 \(x=1\) 处为 \(-1\)，而当 \(x\to\pm\infty\) 时均趋于正无穷，故各有两个不同实交点。两线唯一的公共点 \((1,1)\) 不在抛物线上，所以四点互异；\(k=0\) 时两线重合。


% ITEM: quadratic-example-067
\solutionitem{quadratic-example-067}{例4.67}
过\(P=(0,0)\)作两条不同直线，分别交椭圆\(x^2/9+y^2/4=1\)于\(A,B\)和\(C,D\)。在两斜率均存在时，证明四点共圆当且仅当两斜率互为相反数。分别用推广圆幂与曲线束比较二次项处理。
\solutionanswer
先用推广圆幂。设两单位方向为\(d_1,d_2\)，因两弦过中心，\(|PA||PB|=1/q(d_1)\)、\(|PC||PD|=1/q(d_2)\)。四点共圆必有两乘积相等；反过来，作过\(A,B,C\)的圆，若两乘积相等，圆幂关系便迫使其在\(CD\)上的另一交点为\(D\)。所以四点共圆当且仅当\(q(d_1)=q(d_2)\)。斜率存在时
\[
 q(d_i)=\frac{1/9+k_i^2/4}{1+k_i^2},
\]
两式相等化为\(k_1^2=k_2^2\)。两直线不同，故充要条件为\(k_1+k_2=0\)。

再用曲线束。两直线积为\(L=(y-k_1x)(y-k_2x)\)，椭圆式为\(F=x^2/9+y^2/4-1\)。若四点共圆，两条弦 \(AB,CD\) 的中点都是原点，它们的中垂线交于原点，故圆心为原点。把圆式 \(C\) 的常数项也定为 \(-1\)，则 \(C-F\) 是齐次二次式。它在两条割线上各有非零零点，由齐次性在这两条直线上恒为零，因而是 \(L\) 的倍数。这样，任何经过四点的圆确实属于这个曲线束。要使\(\lambda F+\mu L\)成为圆，交叉项系数\(-\mu(k_1+k_2)\)必须为零，且\(\mu\ne0\)，故仍须\(k_1+k_2=0\)。设\(k_1=k,k_2=-k\)，取\(\mu=1,\lambda=-36(1+k^2)/5\)，两平方项系数均为\(-(4+9k^2)/5\)，确得到经过四点的实圆。


% ITEM: quadratic-exercise-39-01
\solutionitem{quadratic-exercise-39-01}{习题4.39（1）}
抛物线\(y^2=4x\)与\(y=\pm(x-1)\)交于四点。不求交点坐标，求经过它们的圆。
\par
\solutionanswer
将抛物线式记为\(F=y^2-4x\)，两直线积记为\(L=y^2-(x-1)^2\)。取\(2F-L=0\)便得圆\(x^2+y^2-10x+1=0\)，它经过四个共同点。


% ITEM: quadratic-exercise-39-02
\solutionitem{quadratic-exercise-39-02}{习题4.39（2）}
双曲线\(x^2/4-y^2=1\)与\(y=\pm x/3\)交于四点。用曲线束比较二次项，求经过它们的圆。
\par
\solutionanswer
记\(F=x^2/4-y^2-1\)、\(L=y^2-x^2/9\)。取\(8F+9L=x^2+y^2-8=0\)，即所求圆。


% ITEM: quadratic-example-068
\solutionitem{quadratic-example-068}{例4.68}
求经过圆\(x^2+y^2=5\)与直线\(x+y=1\)两交点、且经过原点的圆。再研究经过这两个交点的圆系的圆心轨迹。
\solutionanswer
圆系为\(x^2+y^2-5+\lambda(x+y-1)=0\)。过原点取\(\lambda=-5\)，得\(x^2+y^2-5x-5y=0\)。圆心\((-\lambda/2,-\lambda/2)\)遍历直线\(y=x\)，且每个成员半径平方\(\lambda^2/2+\lambda+5>0\)，故无遗漏。


% ITEM: quadratic-exercise-40-01
\solutionitem{quadratic-exercise-40-01}{习题4.40（1）}
对\(F=x^2+y^2-5\)、\(G=x^2+4y^2-8\)，求\(F+tG=0\)中全部实直线对成员。只令\(x^2\)或\(y^2\)的系数为零，能否找全？
\par
\solutionanswer
成员为\((1+t)x^2+(1+4t)y^2-(5+8t)=0\)。退化参数\(t=-1,-1/4,-5/8\)分别给出\(y=\pm1,x=\pm2,x=\pm2y\)，均为实直线对。只消去平方项系数会漏掉常数项消失时的第三组。


% ITEM: quadratic-exercise-40-02
\solutionitem{quadratic-exercise-40-02}{习题4.40（2）}
求经过圆\(x^2+y^2=4\)与直线\(x=1\)两交点、且圆心在\(y=x\)上的圆。
\par
\solutionanswer
圆系为\(x^2+y^2-4+\lambda(x-1)=0\)，圆心\((-\lambda/2,0)\)。圆心落在\(y=x\)上迫使\(\lambda=0\)，所以所求即\(x^2+y^2=4\)。


% ITEM: quadratic-exercise-40-03
\solutionitem{quadratic-exercise-40-03}{习题4.40（3）}
求圆\(x^2+y^2-2x=0\)、\(x^2+y^2-2y=0\)的根轴与共同点，再求经过共同点及\((2,1)\)的圆。
\par
\solutionanswer
两圆方程相减得根轴 \(x-y=0\)，代回得共同点 \((0,0),(1,1)\)。保留这两个点的圆系可写成 \(x^2+y^2-2x+\lambda(x-y)=0\)。代入 \((2,1)\) 得 \(1+\lambda=0\)，所求圆为 \(x^2+y^2-3x+y=0\)。


% ITEM: quadratic-exercise-40-04
\solutionitem{quadratic-exercise-40-04}{习题4.40（4）}
求\(x^2+y^2=5\)与\(x^2+y^2+xy+x=5\)的全部实交点。选择合适的曲线束成员，说明为什么不必先将两条曲线分别化为标准方程。
\par
\solutionanswer
相减得\(x(y+1)=0\)，故只需分别与\(x=0\)、\(y=-1\)联立。共同点为\((0,\pm\sqrt5),(\pm2,-1)\)。寻找包含交点的退化成员，比逐条曲线化简更直接。


% ITEM: quadratic-exercise-40-05
\solutionitem{quadratic-exercise-40-05}{习题4.40（5）}
（据2012年辽宁卷文科第20题第（1）问）动圆\(x^2+y^2=t^2\)（\(1<t<3\)）与椭圆\(x^2/9+y^2=1\)交于四点。求这四点所成矩形面积的最大值及对应的\(t\)。
\par
\solutionanswer
取第一象限交点 \((x,y)\)，另外三点由对称性得到，矩形面积为 \(4xy\)。只用椭圆条件便有 \(2(x/3)y\le (x/3)^2+y^2=1\)，故面积不超过 \(6\)。取等时 \(x/3=y=1/\sqrt2\)，其所在圆的半径为 \(t=\sqrt{x^2+y^2}=\sqrt5\)，确在 \((1,3)\) 内。因此最大面积就是 \(6\)。


% ITEM: quadratic-example-069
\solutionitem{quadratic-example-069}{例4.69 *}
对\(F=x^2+y^2-5\)、\(G=x^2+4y^2-8\)，写出完整矩阵\(M_F+tM_G\)，用秩判据找出全部退化成员，再将方程分解为直线。说明第三个退化成员与前面两组水平、竖直直线的关系。
\solutionanswer
矩阵\(\operatorname{diag}(1+t,1+4t,-5-8t)\)在\(t=-1,-1/4,-5/8\)退化，秩均为\(2\)。分别给出\(y=\pm1,x=\pm2,x=\pm2y\)。第三组是穿过四个基点的两条对角连接线，前两组则分别按水平、竖直方向配对基点。


% ITEM: quadratic-exercise-41-01
\solutionitem{quadratic-exercise-41-01}{习题4.41（1） *}
求\(x^2+y^2-1+t(x^2+2y^2-3)=0\)的完整矩阵，找出全部退化成员并判断实图形。
\par
\solutionanswer
完整矩阵为\(\operatorname{diag}(1+t,1+2t,-1-3t)\)。退化参数为\(-1,-1/2,-1/3\)，矩阵秩均为\(2\)，实图形依次为\(y=\pm\sqrt2\)、空集、单点\((0,0)\)。中间成员是实数范围无点的复平行直线对。


% ITEM: quadratic-exercise-41-02
\solutionitem{quadratic-exercise-41-02}{习题4.41（2） *}
对两圆\(x^2+y^2-1=0\)、\(x^2+y^2-2x=0\)，写出平方项消去时的齐次二次式，说明它为什么仍由秩小于\(3\)的矩阵表示。其\(z=1\)部分是什么？
\par
\solutionanswer
平方项消去后的齐次二次式为\(2xz-z^2=z(2x-z)\)，矩阵为\(\begin{bmatrix}0&0&1\\0&0&0\\1&0&-1\end{bmatrix}\)，秩\(2\)。\(z=1\)部分为直线\(2x=1\)，另一齐次因子\(z=0\)记录有限平面之外的直线。


% ITEM: quadratic-exercise-41-03
\solutionitem{quadratic-exercise-41-03}{习题4.41（3） *}
设\(M_1,M_2\)线性无关。证明用它们的两个线性无关组合重新作为生成矩阵，所得完整曲线束不变。
\par
\solutionanswer
两个新生成矩阵由一个可逆二阶系数矩阵作用于\((M_1,M_2)\)得到。其逆矩阵又将旧生成元写成新生成元的组合，因此两者张成完全相同的矩阵空间，非零参数对也一一对应，完整曲线束不变。


% ITEM: quadratic-exercise-41-04
\solutionitem{quadratic-exercise-41-04}{习题4.41（4） *}
对\(F=x^2-y^2\)、\(G=xy\)，证明曲线束的全部完整矩阵均退化。写出实图形，说明没有公共直线因子为什么仍可能出现这种情形。
\par
\solutionanswer
完整矩阵末行、末列均为零，故秩至多\(2\)。任意非零成员\(\lambda(x^2-y^2)+\mu xy=0\)的判别式为\(\mu^2+4\lambda^2>0\)，始终是两条过原点的不同实直线。生成曲线没有共同直线因子，但所有成员的二次齐次性使它们都以原点为交会中心，仍可全部退化。


% ITEM: quadratic-exercise-41-05
\solutionitem{quadratic-exercise-41-05}{习题4.41（5） *}
在式\textup{(4.132)}中，设\(M_1\)可逆。证明若某个实数\(t\ne0\)使\(M_1+tM_2\)不可逆，则\(-1/t\)是\(M_1^{-1}M_2\)的特征值。这个关系如何帮助寻找退化成员？
\par

\par\smallskip\noindent 正文所引关系为：
\begin{equation*}
Q_{\lambda,\mu}(Z)=Z^T(\lambda M_1+\mu M_2)Z,
    \qquad M_{\lambda,\mu}=\lambda M_1+\mu M_2 .
\end{equation*}
\solutionanswer
若\((M_1+tM_2)v=0\)，\(v\ne0\)，左乘\(M_1^{-1}\)即得\(M_1^{-1}M_2v=-v/t\)。反向也成立，因此非零实特征值\(\lambda\)对应退化参数\(t=-1/\lambda\)。研究完整束时还应另检验只含\(M_2\)的成员。


% ITEM: quadratic-exercise-41-06
\solutionitem{quadratic-exercise-41-06}{习题4.41（6） *}
对\(F=(x-1)^2+(y+1)^2-1\)、\(G=(x-1)^2+2(y+1)^2-3\)，写出完整矩阵\(M_F+tM_G\)。先平移坐标寻找退化参数，再用原矩阵的秩核对，并解释为何参数不随平移改变。
\par
\solutionanswer
原矩阵为\(\begin{bmatrix}1+t&0&-1-t\\0&1+2t&1+2t\\-1-t&1+2t&1\end{bmatrix}\)。平移\(u=x-1,v=y+1\)后成为\(\operatorname{diag}(1+t,1+2t,-1-3t)\)，退化参数仍为\(-1,-1/2,-1/3\)，秩均为\(2\)。图形依次为\(y+1=\pm\sqrt2\)、空集、点\((1,-1)\)。平移对应同一个可逆合同变换，对所有参数都保持秩，所以不会改变退化参数。


% ITEM: quadratic-exercise-41-07
\solutionitem{quadratic-exercise-41-07}{习题4.41（7） *}
设\(F=x^2+y^2-1\)、\(G=x^2\)。求完整曲线束\(\lambda F+\mu G=0\)的全部退化成员，分别判断其秩与实图形。若只使用\(F+tG=0\)，会遗漏哪个退化成员？
\par
\solutionanswer
矩阵为\(\operatorname{diag}(\lambda+\mu,\lambda,-\lambda)\)。退化成员为：\(\lambda=0\)时\(x^2=0\)，秩\(1\)，重直线；\(\mu=-\lambda\ne0\)时\(y^2=1\)，秩\(2\)，两平行实直线。单参数\(F+tG\)只能取\(\lambda=1\)，会遗漏\(G=x^2\)这一成员。


% ITEM: quadratic-exercise-41-08
\solutionitem{quadratic-exercise-41-08}{习题4.41（8） *}
设三个一次齐次式为\(L=a^TZ\)、\(L_1=b^TZ\)、\(L_2=c^TZ\)。写出\(F=LL_1\)、\(G=LL_2\)的对称矩阵，证明\(\lambda F+\mu G\)的矩阵秩不超过\(2\)，并联系公共直线因子解释。
\par
\solutionanswer
两矩阵为\((ab^T+ba^T)/2\)、\((ac^T+ca^T)/2\)。组合后令\(d=\lambda b+\mu c\)，矩阵为\((ad^T+da^T)/2\)，其像包含在\(\operatorname{span}(a,d)\)，故秩不超过\(2\)。方程本身就是\(L(\lambda L_1+\mu L_2)=0\)，所有成员都含公共直线\(L=0\)。


% ITEM: quadratic-exercise-41-09
\solutionitem{quadratic-exercise-41-09}{习题4.41（9） *}
若曲线束有一个正定的完整矩阵成员，证明生成曲线没有共同实点。反过来，没有共同实点是否必有这样的成员？考察\(F=x^2-y^2-1\)、\(G=x^2-y^2-2\)，分别从方程与矩阵作出判断。
\par
\solutionanswer
共同实点的齐次列\(Z=(x,y,1)^T\ne0\)使每个组合均取零，与某成员正定矛盾。逆命题不成立：所给两方程要求\(x^2-y^2\)同时为\(1,2\)，无共同点；但任意完整矩阵为\(\operatorname{diag}(\lambda+\mu,-\lambda-\mu,-\lambda-2\mu)\)，前两项总是异号或同时为零，不可能正定。
\par
